\section{胶子准PDF算符}\label{sec:operators}

现在可以构造合适的胶子准PDF算符了。
为此，需要把 $J_1^{\mu\nu}$ 的一个指标取为 $z$ 或 $t$，让另一个指标或者跑遍所有洛伦兹分量、或者只跑遍横向分量。在此阶段值得指出：积掉“重夸克”场时，算符 $J_3^{\mu\nu}$ 只产生接触项。这是因为 EOM 算符作用在“重夸克”传播子上会产生一个 $\delta$ 函数。接触项只在 $z_2=z_1$ 时才不为零，这表明当两个 LCOs 缩并到同一时空点时需要额外的重正化。对于非定域胶子准PDF算符的重正化而言，$J_3^{\mu\nu}$ 无关紧要，可以忽略。

由上述讨论，我们确定出构造胶子准PDF算符的如下构件：$J_{1,R}^{zi}$、$J_{1,R}^{ti}$、$J_{1,R}^{z\mu}$。可乘法重正化的胶子非极化准PDF算符可由下式构造：
\begin{align}
 {\cal O}^{1}_R(z_2,z_1) &\equiv  J_{1,R}^{ti}(z_2) \overline J_{1,R}^{ti}(z_1), \non\\
 {\cal O}^{2}_R(z_2,z_1) &\equiv  J_{1,R}^{zi}(z_2)\overline J_{1,R}^{zi}(z_1), \nonumber\\
 {\cal O}^{3}_R(z_2, z_1) &\equiv    J_{1,R}^{ti}(z_2) \overline J_{1,R}^{zi}(z_1),\non\\
 {\cal O}^{4}_R(z_2, z_1)  &\equiv  J_{1,R}^{z\mu}(z_2) \overline J_{1,R,\mu}^{z}(z_1).
\end{align}
积分掉辅助“重夸克”场之后~\cite{Ji:2017oey}，这些算符按如下方式乘法重正化：
\begin{align}
{O}^{1}_R(z_2, z_1) &=Z_{11}^2 e^{\overline{\delta { m}}\Delta z}F^{ti}(z_2) L(z_2,z_1) F^{ti}(z_1),\nonumber\\
 {  O}^{2}_R(z_2,z_1) &=Z_{22}^2 e^{\overline{\delta { m}}\Delta z} F^{zi}(z_2) L(z_2,z_1) F^{zi}(z_1),\nonumber\\
 {O}^{3}_R(z_2,z_1) &= Z_{11}Z_{22} e^{\overline{\delta { m}}\Delta z} F^{ti}(z_2) L(z_2,z_1) F^{zi}(z_1),\nonumber\\
O^{4}_R(z_2,z_1)
& = Z_{22}^2 e^{\overline{\delta { m}}\Delta z} F^{z\mu}(z_2) L(z_2,z_1) F^{z}_{\;\;\mu}(z_1),
 \end{align}
其中 $\Delta z=|z_2-z_1|$。

在大动量极限下，
算符 $O^{i}_R(i=1,2,3,4)$ 彼此之间只相差幂次修正。因此它们属于定义胶子准PDF的同一普适类。
给定以上四个算符，可以用它们的任意组合来研究胶子准PDF，然而这样的组合通常不是可乘法重正化的。一个著名的例子是
\begin{align}
{  O}^{5}_R(z_2,z_1)
 %\equiv   J_{1,R}^{t\mu}(z_2) \overline J^{t}_{1,R,\mu}(z_1)\non\\
 =    -  {  O}^{1}_R(z_2,z_1)-  {  O}^2_R(z_2, z_1)-{  O}^4_R(z_2, z_1),
\end{align}
（去掉迹项后的）它已被用于近期的格点模拟~\cite{Fan:2018dxu}。
但由于 ${  O}^{1}_R(z_2,z_1)$ 与 ${  O}_R^{2,4}(z_2, z_1)$ 的重正化不同，${  O}_R^{5}(z_2,z_1)$ 不是可乘法重正化的。

在以上讨论中我们只考虑了胶子准PDF。如果把胶子准PDF算符插入夸克态，那么只要所有子发散都已被重正化，它就给出紫外有限的贡献。原因在于：准PDF定义在类空间隔上，因此除指数型的质量重正化外不存在其他非定域紫外发散。这表明在上述意义上，胶子准PDF在重正化下不与夸克准PDF混合，混合发生在因子化阶段。这一点已被文献~\cite{Wang:2017qyg} 的单圈计算所证实：胶子准PDF算符的夸克矩阵元紫外有限，但含有与相应劈裂核相联系的 $\ln P_z^2/p^2$（$p^2$ 为红外调节子）项。这一项将与夸克 PDF 匹配，以保证准PDF与 PDF 之间红外发散的相消。
